\section{Hardware Roles: LAN8651, LAN8660, LAN8661}
\label{sec:hardware-roles}

These three parts are not interchangeable. Each has one job.

\subsection{LAN8651 --- host-side MAC-PHY}

\begin{diagram}
    MCU / host
       |
      SPI
       |
    LAN8651
       |
    10BASE-T1S
\end{diagram}

LAN8651 places a programmable MCU onto 10BASE-T1S: a MAC-PHY with a 4-wire SPI host interface implementing the OPEN Alliance TC6 10BASE-T1x MAC-PHY Serial Interface specification (SPI clock up to 25\,MHz), PLCA-capable, with an integrated LDO for 3.3\,V-only designs. It is used in exactly two places in this program: behind the \textbf{Clicker~4} (temporary central node) and behind the \textbf{Gateway MCU} (Section~\ref{sec:arch-b}). \textbf{Preferred initial development hardware: MikroE Two-Wire ETH Click (LAN8651-based, mikroBUS-compatible)} -- it plugs directly into the Clicker~4, and the program does not spend its first PCB revision reproducing functionality that already exists on a known-good commercial board.

\subsection{LAN8660 --- Control Endpoint}

\begin{diagram}
    10BASE-T1S
        |
     LAN8660
        |
    peripheral interface
        |
    motor driver / actuator
\end{diagram}

LAN8660 is Alfred's Control E2B: a fixed-function, MCU-less 10BASE-T1S endpoint (32-pin, 5$\times$5\,mm VQFN) exposing digital I/O through I$^2$C/SPI/UART on its peripheral-facing side. Configuration, control, and diagnostics happen entirely over the network via the standards-based Remote Control Protocol (RCP), orchestrated by a centralized tool (Microchip MPLAB Network Creator) -- there is no Alfred-authored firmware running on the part itself. AEC-Q100 Grade~1 ($-40$ to $+125^\circ$C), ISO~26262 ASIL~B silicon, per Microchip's current (September 2026) public sell sheet (DS00006257C).

\textbf{What is not yet known, and why:} Microchip's public documentation for LAN8660 is a product-brief-level sell sheet, not a pin-exact datasheet -- full register maps, reference schematics, and exact GPIO/analog capability are gated behind Microchip's secure/NDA documentation program. Requesting that access is a Week~1, critical-path task (Section~\ref{sec:timeline}), and the exact motor-driver/actuator interface (Section~\ref{sec:lan8660-validation}) is deliberately left open until that documentation, or evaluation hardware, is in hand -- consistent with the instruction not to lock an actuator choice before the real interface is understood.

\subsection{LAN8661 --- Lighting Endpoint}

\begin{diagram}
    10BASE-T1S
        |
     LAN8661
        |
    lighting-driver interface
        |
       LEDs
\end{diagram}

Same silicon family and MCU-less architecture as LAN8660, specialized for lighting: it bridges T1S video/remote-control data to an external LED driver IC rather than exposing generic sensor/actuator I/O. Same documentation caveat applies -- the exact LED-driver IC (Section~\ref{sec:lan8661-validation}) is chosen after review of Microchip's LAN8661 reference material, not assumed in advance.

\subsection{Why the roles cannot be blurred}

\begin{center}
\small
\begin{tabularx}{\linewidth}{@{}p{3.1cm} X@{}}
\toprule
\textbf{Attempt} & \textbf{Why it fails} \\
\midrule
Use LAN8660/LAN8661 as the Gateway's network chip & RCP is a fixed, network-managed control-plane protocol for peripheral-style I/O, not a general transport an MCU can shape into arbitrary normalized-CAN-frame payloads; there is no host-side ingress path for a companion MCU to inject application-defined traffic. \\
Put an MCU behind LAN8660/LAN8661 to add custom logic & Wrong interface direction -- their I$^2$C/SPI/UART pins face outward to peripherals, not inward to a host controller. There is nothing for a companion MCU to talk to that would change the endpoint's behavior beyond what RCP already configures. \\
One populated-selectively PCB serving both E2B and Gateway & Populating LAN8660/LAN8661 makes an MCU, its SPI link, and any firmware image irrelevant, since the endpoint needs none of it -- that is not a population variant of the same board, it is a structurally different device sharing a board outline at most (Section~\ref{sec:platform-strategy}). \\
\bottomrule
\end{tabularx}
\end{center}
